% recommendations.tex --- evidence-gated research and control agenda

\section{Research and Control Recommendations}\label{sec:recommendations}

The present evidence supports an offline attribution study and a threat-model
agenda.  It does not support transaction blocking, fraud claims about
registry-associated owners, or a production rollout.  Recommendations are
therefore organized as evidence gates and control hypotheses rather than as a
deployment roadmap.

\subsection{Keep the Programs Separate}

Institutions and researchers should maintain separate objectives, labels, and
success criteria for:
\begin{enumerate}[leftmargin=*,nosep]
  \item \textbf{owner or automation attribution}, including registry
    association and verified controller type;
  \item \textbf{control-plane security}, including prompt or tool injection,
    credential compromise, mandate escalation, and provenance loss;
  \item \textbf{fraud}, using authorization, deception, victim, loss, and case
    outcomes;
  \item \textbf{AML and sanctions}, using entity, source-of-funds, typology, and
    investigative evidence; and
  \item \textbf{market abuse}, using trades, beneficial-control hypotheses,
    venue mechanics, and market impact.
\end{enumerate}

An attribution score may later be tested as one contextual feature in the other
programs.  It should not supply their ground truth or trigger an adverse action
by itself.

\subsection{Control Hypotheses by Data Plane}

The causal model points first to preventive and evidentiary controls that bind
delegated authority.  \Cref{tab:control_hypotheses} lists candidate controls;
the present experiment does not validate their operational effectiveness.

\begin{table}[t]
  \centering
  \caption{Candidate controls implied by the causal model.}
  \label{tab:control_hypotheses}
  \small
  \begin{tabularx}{\textwidth}{@{}p{2.6cm} p{5.0cm} X@{}}
    \toprule
    \textbf{Plane} & \textbf{Candidate control} & \textbf{Evidence to retain} \\
    \midrule
    Control & Authenticated, bounded mandate; tool least privilege; separation
      of content from executable instructions; versioned policy & Principal,
      mandate, prompt/message provenance, tool call and policy decision \\
    Wallet authorization & Per-mandate amount and count limits; beneficiary and
      contract scopes; simulation; step-up; revocation and emergency stop &
      Custody context, evaluated rule, approval, signature and denial log \\
    Settlement & Typology-specific monitoring of new beneficiaries, shared
      funding, flow and entity graph; rail-level rate limits where justified &
      Transaction, fee, contract, chain, counterparty and entity-resolution data \\
    Outcome & Delivery confirmation, dispute and loss capture, investigator
      adjudication, model-feedback controls & Service result, customer claim,
      case rationale, disposition and intervention outcome \\
    \bottomrule
  \end{tabularx}
\end{table}

Mandate-bearing payment protocols and programmable wallets provide concrete
design patterns for the first two planes
\cite{ap2protocol2025,li2026a402,metamask2026agentwallet}.  Their presence does
not prove that a given deployment is safe; it makes the policy and evidence
requirements testable.

\subsection{Evidence Gates Before Operational Use}

The following gates should be satisfied in order.  Failure at an earlier gate
invalidates later performance claims.

\begin{enumerate}[leftmargin=*,nosep,label=\textbf{G\arabic*.}]
  \item \textbf{Construct gate.} State the target, unit of analysis, role,
    decision, and excluded interpretations.  Rename the current target
    registered-identity owner discrimination.
  \item \textbf{Label gate.} Verify registry timing and owner semantics; classify
    contracts and service entities; retain an unknown class; obtain independent
    outcome labels for any fraud, AML, or market-abuse experiment.
  \item \textbf{Split and leakage gate.} Use post-registration observations,
    temporal holdouts, and controller- or entity-grouped splits so linked
    addresses cannot appear across training and test partitions.
  \item \textbf{Baseline gate.} Compare with registry lookup, simple rules,
    calibrated supervised models, and graph/entity baselines.  Report
    incremental lift rather than performance in isolation.
  \item \textbf{Adversarial gate.} Use independently produced, held-out attack
    generators, stochastic and adaptive policies, realistic address histories,
    and legitimate automated twins.
  \item \textbf{Operational gate.} Evaluate prospectively in shadow mode using
    calibration, precision--recall, precision at investigator capacity, false
    positives by entity type, intervention outcomes, and stability under drift.
  \item \textbf{Governance gate.} Complete legal, privacy, model-risk, security,
    and human-review assessments for the actual institution and jurisdiction.
\end{enumerate}

The current study does not pass the label gate for fraud or the operational
gate.  Its real-data result should remain a research benchmark until those gaps
are closed.

\subsection{Evaluation and Falsification Criteria}

Future work should preregister both hypotheses and stop conditions.  At minimum,
the research claim should be rejected or narrowed if:
\begin{itemize}[leftmargin=*,nosep]
  \item registry-associated separation disappears under post-registration,
    entity-grouped temporal evaluation;
  \item the score adds no lift beyond registry status and standard graph or
    transaction baselines;
  \item performance cannot meet a prespecified false-positive or investigator
    capacity constraint;
  \item held-out adaptive generators evade the tested signatures; or
  \item joined mandate and outcome evidence shows that the flagged behavior was
    authorized and legitimate.
\end{itemize}

For rare financial-crime outcomes, ROC-AUC alone is insufficient.  Studies
should report class prevalence, precision--recall curves, calibration,
confidence intervals, alert volume, case-level rather than address-level
metrics, and the effect of any intervention.  Synthetic injections should be
described as implementation or design-envelope checks unless external attack
generation and realistic histories justify a stronger term.

\subsection{Privacy, Retention, and Accountability}\label{subsec:privacy}

Public availability does not by itself make on-chain data anonymous or exempt
from data-protection law.  Addresses and linked behavioral profiles can become
personal data when they are reasonably associated with a natural person;
pseudonymization reduces some risks but does not remove all obligations
\cite{gdpr2016,ccpa2018}.  The applicable lawful basis, purpose limitation,
access controls, retention period, data-subject rights, and impact assessment
must be determined for the actual use case and jurisdiction.

Cross-platform linkage and control-plane transcripts are particularly
sensitive because they can join identity, communications, and financial
behavior.  They should not be collected ``just in case'' or retained under a
blanket period.  Collection and retention should be tied to a documented
purpose, minimized, access-controlled, and reviewed by qualified privacy and
legal functions.  Financial-crime information sharing and reporting likewise
requires scheme- and jurisdiction-specific authority; the present technical
study supplies no legal safe harbor.

Finally, accountability should follow the role and mandate chain.  The relevant
questions are who authorized the agent, which policy and wallet enforced the
delegation, what evidence was preserved, and what outcome occurred.  Treating
software as the accountable ``fraudster'' obscures those questions and weakens
both prevention and investigation.
